\chapter{Complex Numbers, Phasors, and Impedance}
\section{Why Complex Numbers Appear}
A sinusoid can be represented as the real part of a complex exponential:
\[
x(t)=\sqrt{2}\,\Re\{X\ee^{\jj\omega t}\}.
\]
The factor \(\sqrt{2}\) appears because this book uses RMS phasor magnitudes by
default. Differentiation becomes multiplication by \(\jj\omega\):
\[
\frac{\dd}{\dd t}\left(X\ee^{\jj\omega t}\right)
=\jj\omega X\ee^{\jj\omega t}.
\]
That one fact turns differential equations into algebra for sinusoidal
steady-state analysis.

\section{Rectangular and Polar Form}
\[
Z=R+\jj X=|Z|\angle\theta,
\qquad
|Z|=\sqrt{R^2+X^2},
\qquad
\theta=\operatorname{atan2}(X,R).
\]
Rectangular form is natural for adding series impedances. Polar form is natural
for multiplication, division, and phase interpretation.

\begin{mistakebox}
\(\tan^{-1}(X/R)\) alone can put the angle in the wrong quadrant. Use
\(\operatorname{atan2}(X,R)\) or reason from the signs of both coordinates.
\end{mistakebox}

\section{Impedance of R, L, and C}
\[
Z_R=R,\qquad
Z_L=\jj\omega L,\qquad
Z_C=\frac{1}{\jj\omega C}=-\jj\frac{1}{\omega C}.
\]
An inductor has positive reactance. A capacitor has negative reactance. At low
frequency, a capacitor's impedance looks large and an inductor's looks small. At
high frequency, the roles reverse.

\begin{figure}[htbp]
\centering
\begin{circuitikz}
\draw (0,0) to[sV,l=\(v_s\)] (0,2)
      to[R,l=\(R\)] (2.6,2)
      to[L,l=\(L\)] (5.2,2)
      to[C,l=\(C\)] (5.2,0)
      -- (0,0);
\end{circuitikz}
\caption{A series RLC circuit whose impedance is \(R+\jj(\omega L-1/\omega C)\).}
\end{figure}

\section{Phasor Lead and Lag}
For a capacitor, current leads voltage by \(\SI{90}{\degree}\). For an
inductor, voltage leads current by \(\SI{90}{\degree}\). A positive impedance
angle means inductive behavior; a negative impedance angle means capacitive
behavior.

\begin{exambox}
If an exam question says ``lead'' or ``lag,'' sketch a tiny phasor diagram before
answering. Most mistakes come from remembering the right \(\SI{90}{\degree}\)
relationship with the direction reversed.
\end{exambox}

\section{Admittance}
Admittance is the reciprocal of impedance:
\[
Y=\frac{1}{Z}=G+\jj B.
\]
It is especially convenient for parallel networks, because admittances add:
\[
Y_{\mathrm{eq}}=\sum_k Y_k.
\]

\section{Complex Power}
With RMS phasors and passive sign convention, complex power is
\[
S=VI^*=P+\jj Q.
\]
The complex conjugate matters. Real power \(P\) is average power; reactive power
\(Q\) represents energy sloshing between fields and sources.

\section{Worked Example}
Find the impedance of a \(\SI{50}{\ohm}\) resistor in series with
\(\SI{10}{\micro\henry}\) at \(\SI{7.1}{\mega\hertz}\).
\[
\omega=2\pi(7.1\times10^6)=\SI{4.46e7}{\radian\per\second},
\qquad
X_L=\omega L=446~\Omega.
\]
Thus
\[
Z=50+\jj 446~\Omega,
\qquad
|Z|=\SI{449}{\ohm},
\qquad
\theta\approx \SI{83.6}{\degree}.
\]
The circuit is strongly inductive.
